problem:  
Let \((M_i^n,g_i,p_i)\) be a non-collapsing sequence of complete Riemannian manifolds with \(\operatorname{Ric}_{g_i}\geq -(n-1)\) converging in the pointed Gromov–Hausdorff sense to a limit space \((X,d,p)\). By Cheeger–Colding theory, each tangent cone at the basepoint \(p\in X\) is a metric cone \(C(Z)\). Denote by \(k(p)\) the maximal integer for which the tangent cone splits off an isometric \(\mathbb{R}^{k(p)}\) factor.

By results of Kapovitch–Wilking, for sufficiently large \(i\), the fundamental groups \(\pi_1(M_i,p_i)\) contain nilpotent subgroups of uniformly bounded index and step. However, current results do not directly bound the *rank* of these nilpotent subgroups in terms of the geometric structure of the limit space.

**Conjectural problem:**  
Does the Euclidean splitting dimension \(k(p)\) of the tangent cone at \(p\in X\) give a sharp quantitative upper bound on the nilpotent rank of \(\pi_1(M_i,p_i)\)? In particular, is it true that for all sufficiently large \(i\),
\[
\operatorname{rank}_{\mathrm{nil}}\big(\pi_1(M_i,p_i)\big) \le k(p)?
\]
Moreover, can the rank drop under pointed measured Gromov–Hausdorff convergence only when the number of Euclidean factors in the tangent cone decreases?

Why is it a "good" problem:  
This question isolates the *quantitative geometric–algebraic correspondence* that remains unproven at the interface of Cheeger–Colding limit theory and the Kapovitch–Wilking structural theory of π₁. The virtual nilpotency of π₁ is qualitatively known, but no theorem yet ties the numerical invariant “nilpotent rank” directly to the *metric splitting dimension* of the Ricci-limit space. If such a bound existed—and especially if it were sharp—it would create a precise dictionary linking analytic splitting phenomena (Euclidean factors in tangent cones) with the algebraic stratification of π₁(M_i).  

The problem is conceptually simple (“Does the tangent cone’s Euclidean dimension bound the algebraic rank of the fundamental group?”) but profoundly deep, intertwining metric geometry, Ricci curvature analysis, and asymptotic group theory. It targets the missing quantitative parameter unifying geometric convergence and algebraic stability, pushing beyond existing theorems and potentially leading to a new classification principle for π₁ in Ricci-limit geometry.